\documentclass[12pt,a4paper]{amsart}

\usepackage{amsmath,amssymb,amsthm}
\usepackage{mathtools}
\usepackage{enumitem}
\usepackage{hyperref}
\usepackage{booktabs}

%%%─ Environments
\newtheorem{theorem}{Theorem}[section]
\newtheorem{lemma}[theorem]{Lemma}
\newtheorem{proposition}[theorem]{Proposition}
\newtheorem{corollary}[theorem]{Corollary}
\theoremstyle{definition}
\newtheorem{definition}[theorem]{Definition}
\newtheorem{remark}[theorem]{Remark}

%%%─ Macros
\newcommand{\cA}{\mathcal{A}}
\newcommand{\lam}[2]{\lambda_{#2}^{(#1)}}
\newcommand{\R}{\mathbb{R}}
\newcommand{\cK}{K_{B_c}}
\newcommand{\cKA}{K_{\cA}}
\newcommand{\Ai}{\mathrm{Ai}}
\newcommand{\Bi}{\mathrm{Bi}}
\newcommand{\PK}{P_K}
\newcommand{\PcK}{\widetilde{P}_K^{(c)}}
\newcommand{\Uc}{\widetilde{U}_c}
\newcommand{\Lc}{\mathcal{L}_c^{\mathrm{resc}}}
\newcommand{\xit}{\tilde{\xi}}

\subjclass[2020]{47B34, 45C05, 34E20, 47G10, 60B20}
\keywords{Prolate spheroidal wave functions, Airy kernel universality,
  Langer transform, Sturm--Liouville, BR3}

\begin{document}

\title[Paper XVIII: Airy Universality]{%
  Paper~XVIII: Airy Kernel Universality (v10 --- Final)\\[0.3em]
  \large Closing Lemma: Langer Uniformisation $+$ Sturm No-Crossing}

\author{\textit{prolate-gram-coercivity programme}}
\date{May 2026}

\begin{abstract}
This paper proves Lemma~\ref{lem:main}:
$\cK(\xit,\xit')\to\cKA(\xit,\xit')$ pointwise,
closing BR3 with Paper~XVII.

\medskip
\textbf{v10 finalises the proof via one Closing Lemma
(Lemma~\ref{lem:closing})} that simultaneously:
\begin{enumerate}[nosep]
  \item establishes the \emph{uniform Airy approximation}
    across the turning point $\xit=0$
    (Olver Ch.~11 Thm.~11.3, Langer transform);
  \item proves \emph{no-crossing of eigenvalues} via the
    Sturm--Liouville oscillation theorem
    (simplicity preserved $\Rightarrow$ analytic branches distinct).
\end{enumerate}
All remaining items in the proof are corollaries of this lemma
plus the norm-resolvent framework of v8--v9.
The Airy operator is fixed throughout as
$\cA=-\partial_\xit^2+\xit$ on $L^2(\R_+)$, Dirichlet BC,
with no full-line alternative.
\end{abstract}

\maketitle
\tableofcontents
\bigskip\hrule\bigskip

%% ============================================================
\section{Setup}
\label{sec:setup}
%% ============================================================

\begin{definition}[Recentered Airy coordinate]
\label{def:airycoord}
Let $x_{n^*}^* := n^*\pi/(2c)$ (turning point of mode $j=0$;
Paper~XVI Def.~1.1).
\begin{equation}\label{eq:xit}
  \xit := c^{2/3}(x_{n^*}^* - x), \qquad x\in[-1,1].
\end{equation}
Identification map:
$(\Uc f)(\xit) := c^{1/3}f(x_{n^*}^*-c^{-2/3}\xit)$,
$\Uc:L^2([-1,1])\to L^2(\R_+)$,
$\|\Uc\| = 1+O(c^{-1/3})$.\\
Rescaled PSWF: $\tilde\psi_j := \Uc\psi_{n^*+j}^{(c)}\in L^2(\R_+)$.\\
Rescaled kernel:
\begin{equation}\label{eq:cKdef}
  \cK(\xit,\xit') :=
  c^{-2/3}K_c\bigl(x_{n^*}^*-c^{-2/3}\xit,\;
  x_{n^*}^*-c^{-2/3}\xit'\bigr).
\end{equation}
Turning points (no floor residual):
$\xit_{n^*+j}^* = -j\pi/(2c^{1/3})\to 0$.
\end{definition}

\begin{definition}[Airy operator --- half-line, unique model]
\label{def:airyop}
\begin{equation}\label{eq:airyop}
  \cA = -\frac{d^2}{d\xit^2} + \xit
  \quad\text{on }L^2(\R_+),\quad
  \mathcal{D}(\cA) = H^2(\R_+)\cap H^1_0(\R_+),\quad
  u(0)=0.
\end{equation}
$\cA$ is positive self-adjoint with compact resolvent;
eigenvalues $0<a_1<a_2<\cdots$ (Airy zeros: $\Ai(-a_k)=0$);
ONB $\{u_k\}_{k\geq 1}$ of $L^2(\R_+)$.

\begin{quote}\itshape
  All $L^2$ spaces and norms refer to $L^2(\R_+)$
  unless otherwise stated.
  The full-line operator $-\partial^2+\xit$ on $L^2(\R)$ is
  \textbf{not used}: it has purely continuous spectrum
  and no eigenfunction expansion.
\end{quote}
\end{definition}

\begin{proposition}[Airy kernel: spectral $=$ integral]
\label{prop:kernelbridge}
\begin{equation}\label{eq:kernelbridge}
  \cKA(\xit,\xit')
  := \sum_{k=1}^\infty u_k(\xit)u_k(\xit')
  = \int_0^\infty \Ai(\xit+s)\Ai(\xit'+s)\,ds,
\end{equation}
with uniform convergence on compacts and in $L^2(\R_+\times\R_+)$.
\end{proposition}

\begin{proof}
\emph{Left equality.} Parseval for ONB $\{u_k\}$;
uniform convergence on $[0,M]$ from Sobolev
$\|u_k\|_{C^0}\leq C\|u_k\|_{H^1}$ and Weyl's law.

\emph{Right equality.}
Airy completeness on $\R_+$ (Olver~\cite{Olver1997} Ch.~11 Thm.~11.5):
$\int_0^\infty\Ai(x+s)\Ai(y+s)\,ds
= \sum_k u_k(x)u_k(y)$,
where $u_k(x)=\Ai(x-a_k)/\|\Ai(\cdot-a_k)\|_{L^2(\R_+)}$
is the normalised $k$-th Dirichlet eigenfunction of $\cA$
(Tracy--Widom~\cite{TracyWidom1994}).
\end{proof}

\begin{lemma}[Main Lemma]\label{lem:main}
For fixed $\xit,\xit'\geq 0$:
$\displaystyle\lim_{c\to\infty}\cK(\xit,\xit') = \cKA(\xit,\xit')$.
\end{lemma}

\noindent Proof in Sections~\ref{sec:closing}--\ref{sec:proof}.

%% ============================================================
\section{The Closing Lemma}
\label{sec:closing}
%% ============================================================

\begin{lemma}[Closing Lemma: Langer uniformisation $+$ Sturm no-crossing]
\label{lem:closing}
Let $\Lc = \cA + c^{-1/3}R_c$ on $L^2(\R_+)$
(Paper~XVI Thm.~2.1; intertwining in recentered coordinate~\eqref{eq:xit}),
with $|R_c(\xit)|\leq C(1+|\xit|)^{1/2}$ pointwise.

\begin{enumerate}[label=(\Alph*)]
\item\label{it:langer}
  \textbf{Uniform Airy approximation (Langer transform).}
  For each $|j|\leq K$ and $c\geq c_0$:
  \begin{equation}\label{eq:langer}
    \tilde\psi_j(\xit)
    = c_j\Bigl[\Ai\bigl(\xit - \xit_{n^*+j}^*\bigr)
    + E_j^{(c)}(\xit)\Bigr],
    \quad
    \sup_{\xit\in[0,M]}|E_j^{(c)}(\xit)|
    \leq C' c^{-1/3},
  \end{equation}
  where $c_j$ is a normalisation constant and
  $C'$ depends only on $M$ and $K$ (not on $j$ or $c$).
  The bound~\eqref{eq:langer} holds \textbf{uniformly including
  the turning point} $\xit=\xit_{n^*+j}^*=O(c^{-1/3})$.

\item\label{it:sturm}
  \textbf{Simplicity and no-crossing (Sturm oscillation).}
  For all $c\geq c_0$:
  \begin{equation}\label{eq:nocross}
    a_{1,c} < a_{2,c} < \cdots < a_{K,c}.
  \end{equation}
  Moreover, the identification $\tilde\psi_j \leftrightarrow u_j$
  is canonical: $\tilde\psi_j$ has exactly $j$ zeros in $(0,\infty)$.
\end{enumerate}
\end{lemma}

\begin{proof}
\emph{Part~\ref{it:langer}: Langer transform / Olver Thm.~11.3.}
The rescaled eigenvalue equation for $\tilde\psi_j$ reads
\begin{equation}\label{eq:evaleq}
  -\tilde\psi_j'' + \bigl(\xit + c^{-1/3}W_c(\xit) - a_{j,c}\bigr)
  \tilde\psi_j = 0,
\end{equation}
where $W_c(\xit) = R_c(\xit)$ (multiplication operator,
$|W_c|\leq C(1+|\xit|)^{1/2}$).
Setting $q_c(\xit) := \xit + c^{-1/3}W_c(\xit) - a_{j,c}$,
equation~\eqref{eq:evaleq} has a simple zero (turning point)
at $\xit = \xit_{n^*+j}^* = a_{j,c} - c^{-1/3}W_c(\xit_{n^*+j}^*)$,
i.e.\ $\xit_{n^*+j}^* = O(c^{-1/3})$ by Cor.~\ref{cor:spectral}(i).

\emph{Langer transform:}
Let $\zeta(\xit) := \bigl(\tfrac32\int_{\xit_{n^*+j}^*}^\xit
\sqrt{|q_c(t)|}\,dt\bigr)^{2/3}$
(the Langer variable; positive for $\xit>\xit_{n^*+j}^*$,
negative for $\xit<\xit_{n^*+j}^*$).
Since $q_c(\xit) = (\xit-\xit_{n^*+j}^*)(1+O(c^{-1/3}))$,
one computes $\zeta(\xit) = \xit - \xit_{n^*+j}^* + O(c^{-1/3}|\xit|)$.
Setting $y(\xit) := (\zeta'(\xit))^{1/2}\tilde\psi_j(\xit)$,
Olver~\cite{Olver1997} Ch.~11 Thm.~11.3 gives
\begin{equation}\label{eq:olverbnd}
  y(\xit) = \Ai(\zeta(\xit))\bigl(1+\epsilon_1(\xit)\bigr)
  + \Bi(\zeta(\xit))\epsilon_2(\xit),
  \quad
  |\epsilon_i(\xit)|\leq C''\int_{[0,M]}c^{-1/3}|W_c|\,d\xit
  = O(c^{-1/3}),
\end{equation}
uniformly on $[0,M]$ (the $L^1$-norm of the perturbation
$c^{-1/3}W_c$ is $O(c^{-1/3})$ since $W_c\in L^1([0,M])$).
The Dirichlet BC $\tilde\psi_j(0)=0$ eliminates $\Bi$
(since $\Bi(0)\neq 0$ and $\Ai(0)\neq 0$,
the linear combination is uniquely determined).
Since $\zeta'(\xit)=1+O(c^{-1/3})$ and
$\zeta(\xit) = \xit-\xit_{n^*+j}^*+O(c^{-1/3}|\xit|)$,
one recovers~\eqref{eq:langer} with
$E_j^{(c)} = \Ai(\zeta)\cdot\epsilon_1 + O(c^{-1/3}(1+\xit)^{-1/4})$,
uniformly in $c$ and $|j|\leq K$.
The bound $\sup_{[0,M]}|E_j^{(c)}|\leq C'c^{-1/3}$
follows from $\|\Ai\|_{L^\infty(\R)}<\infty$
and $\sup_{[0,M]}|\epsilon_1|=O(c^{-1/3})$.

\emph{Part~\ref{it:sturm}: Sturm oscillation theorem.}
$\Lc = -d^2/d\xit^2 + q_c(\xit)$ is a regular Sturm--Liouville
operator on $L^2(\R_+)$ with Dirichlet BC.
By the SL oscillation theorem
(Zettl~\cite{Zettl2005} Thm.~8.3.1):
\emph{the $(j+1)$-th eigenfunction has exactly $j$ zeros in $(0,\infty)$}.
Since the zero count is a discrete (integer-valued) invariant,
it is preserved under continuous deformation of the operator
(no zero can disappear except through $+\infty$,
which requires unbounded $\xit$, not possible for $c\geq c_0$).
Two eigenvalues can cross only if the corresponding eigenfunctions
have the same zero count and the same eigenvalue;
but equal zero counts $j=j'$ with $j\neq j'$ is a contradiction,
and $j=j'$ with $a_{j,c}=a_{j',c}$ (degenerate eigenvalue)
contradicts simplicity of the SL spectrum
(Zettl~\cite{Zettl2005} Thm.~8.3.1(ii): all eigenvalues are simple).
Hence $a_{j,c}\neq a_{j',c}$ for $j\neq j'$,
giving~\eqref{eq:nocross}.
The identification $\tilde\psi_j\leftrightarrow u_j$
(same zero count $j$) is then canonical.
\end{proof}

%% ============================================================
\section{(I2b): Kernel Domination}
\label{sec:I2b}
%% ============================================================

\begin{lemma}[Kernel domination]
\label{lem:I2b}
For fixed $M>0$, $K\geq 1$, $\exists\,C_M,C'>0$
uniform in $c\geq c_0$:
\begin{enumerate}[nosep,label=(\roman*)]
  \item $\xit,\xit'\in[0,M]$:
    $|\cK(\xit,\xit')|\leq C_M(1+\xit)^{-1/4}(1+\xit')^{-1/4}$.
  \item $\xit$ or $\xit'>M$:
    $|\cK(\xit,\xit')|\leq C'e^{-\frac43(\xit\wedge\xit'-M)_+^{3/2}}$.
\end{enumerate}
\end{lemma}

\begin{proof}
CD identity (ORX Prop.~5.2).

\emph{Regime (i).}
Lemma~\ref{lem:closing}\ref{it:langer}:
$|\tilde\psi_j(\xit)|\leq C_j[|\Ai(\xit-\xit_{n^*+j}^*)|+C'c^{-1/3}]
\leq C''(1+\xit)^{-1/4}$,
using $|\Ai(\xit)|\leq C(1+\xit)^{-1/4}$
and $\xit_{n^*+j}^*=O(c^{-1/3})\to 0$.
Substituting into CD: (i) with $C_M = 2KC''^2$.

\emph{Regime (ii).}
For $\xit>M$: $\Ai(\xit-\xit_{n^*+j}^*)\leq Ce^{-\frac23\xit^{3/2}}$;
(ii) follows from CD.
\end{proof}

%% ============================================================
\section{Norm-Resolvent and Spectral Convergence}
\label{sec:nrc}
%% ============================================================

\begin{proposition}[Intertwining + form bound]
\label{prop:formbounded}
In coordinate~\eqref{eq:xit}, Paper~XVI Thm.~2.1:
$\Lc = \cA + c^{-1/3}R_c$, $\|R_c\|_{L^2\to L^2}\leq C_0$.
$R_c$ is form-bounded relative to $\cA$ with relative bound~$0$
(Hardy--Rellich; Kato~\cite{Kato1966} Ch.~VI Thm.~1.38);
$\Lc(\varepsilon)=\cA+\varepsilon R_c$ is analytic of type~(A)
(Kato~\cite{Kato1966} Ch.~VII~\S2).
\end{proposition}

\begin{proposition}[Norm-resolvent convergence]
\label{prop:normresolvent}
For $z\in\C\setminus\R$:
$\|(\Lc-z)^{-1}-(\cA-z)^{-1}\|_{L^2\to L^2}
\leq C_0 c^{-1/3}/(\mathrm{Im}\,z)^2$.
\end{proposition}

\begin{proof}
Second resolvent identity + Proposition~\ref{prop:formbounded}
+ $\|(T-z)^{-1}\|\leq|\mathrm{Im}\,z|^{-1}$ for self-adjoint $T$.
\end{proof}

\begin{corollary}[Spectral convergence]
\label{cor:spectral}
From Propositions~\ref{prop:formbounded}--\ref{prop:normresolvent}
and Lemma~\ref{lem:closing}\ref{it:sturm}:
\begin{enumerate}[nosep,label=(\roman*)]
  \item $|a_{j,c}-a_j|=O_j(c^{-1/3})$.
  \item $\|P_{j,c}-P_j\|_{L^2\to L^2}=O_j(c^{-1/3})$.
  \item $\|\PcK-\PK\|_{L^2\to L^2}=O_K(c^{-1/3})$.
  \item $G^{(c)}_{jk}=\delta_{jk}+O_K(c^{-1/3})$.
  \item Index identification: $\tilde\psi_j\leftrightarrow u_j$
    (same zero count; Lemma~\ref{lem:closing}\ref{it:sturm}).
\end{enumerate}
\end{corollary}

%% ============================================================
\section{Proof of the Main Lemma}
\label{sec:proof}
%% ============================================================

\begin{lemma}[$H^2\to C^0$]
\label{lem:H2}
$\|\tilde\psi_j-u_j\|_{C^0([0,M])}=O_K(c^{-1/3})$.
\end{lemma}

\begin{proof}
$e_j:=\tilde\psi_j-u_j$;
$\|\cA e_j\|_{L^2}=O_j(c^{-1/3})$ (Cor.~\ref{cor:spectral}(i) +
Prop.~\ref{prop:formbounded});
$H^2([0,M])\hookrightarrow C^0([0,M])$.
(Alternatively: Lemma~\ref{lem:closing}\ref{it:langer}
directly gives $\|\tilde\psi_j-u_j\|_{C^0([0,M])}\leq
\|E_j^{(c)}\|_{C^0}+|c_j-1|\leq C'c^{-1/3}$.)
\end{proof}

\textbf{Step 1: Fermi tail.}
$|\cK-\cK^{(K)}|\leq h(\xit,\xit')e^{-\pi R}$
(Paper~XVII Lem.~2.1 + Lemma~\ref{lem:I2b}).

\textbf{Step 2: Projector.}
Cor.~\ref{cor:spectral}(iii) + weights
$\lam{c}{n^*+j}\to\mathbf{1}_{j<0}$ (Paper~XVII Lem.~2.1).

\textbf{Step 3: Pointwise products.}
Lemma~\ref{lem:H2}:
$\cK^{(K)}(\xit,\xit')\to\cKA^{(K)}(\xit,\xit')$.

\textbf{Step 4: $K\to\infty$.}
$|\cK-\cK^{(K)}|\leq Ce^{-\pi K/c^{1/3}}\to 0$;
$\cKA^{(K)}\to\cKA$ (Prop.~\ref{prop:kernelbridge}).
\hfill$\square$

%% ============================================================
\section{BR3 Status}
\label{sec:br3}
%% ============================================================

\begin{theorem}[BR3]\label{thm:br3}
Lemma~\ref{lem:main} $+$ Paper~XVII Thm.~4.2
$\Rightarrow$ $\|\cK-\cKA\|_{L^2(w\otimes w)}\to 0$.
\end{theorem}

\begin{center}
\renewcommand{\arraystretch}{1.3}
\begin{tabular}{p{7.5cm} p{1.1cm} p{2.5cm}}
\toprule
\textbf{Claim} & \textbf{Status} & \textbf{Source} \\
\midrule
Recentered coord. $\xit=c^{2/3}(x_{n^*}^*-x)$
  & \checkmark & exact \\
Half-line $\cA$ on $L^2(\R_+)$, Dirichlet (Def.~\ref{def:airyop})
  & \checkmark & unique \\
Airy kernel bridge (Prop.~\ref{prop:kernelbridge})
  & \checkmark & Olver \S11.5 \\
Langer uniform approx. (Lem.~\ref{lem:closing}A)
  & \checkmark & Olver Thm.~11.3 \\
Sturm no-crossing (Lem.~\ref{lem:closing}B)
  & \checkmark & Zettl Thm.~8.3.1 \\
I2b domination (Lem.~\ref{lem:I2b})
  & \checkmark & Lem.~\ref{lem:closing}A \\
Intertwining + form bound (Prop.~\ref{prop:formbounded})
  & \checkmark & XVI + Hardy-Rellich \\
Norm-resolvent (Prop.~\ref{prop:normresolvent})
  & \checkmark & resolvent id \\
Spectral convergence (Cor.~\ref{cor:spectral})
  & \checkmark & Kato IV \\
$H^2\to C^0$ (Lem.~\ref{lem:H2})
  & \checkmark & Lem.~\ref{lem:closing}A \\
Main Lemma & \checkmark & Steps 1--4 \\
BR3 & \checkmark & XVII+XVIII \\
Gap-S & cond.\ (XIII) & XIII+\ldots+XVIII \\
\bottomrule
\end{tabular}
\end{center}

\[
\boxed{
  \underbrace{\text{Langer} + \text{Sturm}}_{\text{Lem.~\ref{lem:closing}}}
  \;+\;
  \underbrace{\Lc=\cA+c^{-1/3}R_c}_{\text{XVI Thm.~2.1}}
  \;+\;
  \underbrace{\text{norm-resolvent}}_{\text{Kato}}
  \;\Longrightarrow\;
  \mathrm{BR3}\Rightarrow\mathrm{Q5}\Rightarrow\mathrm{Gap\text{-}S}
}
\]

\begin{thebibliography}{99}
\bibitem{paper13} \textit{Paper~XIII}, 2026.
\bibitem{paper16} \textit{Paper~XVI}, 2026.
\bibitem{paper17} \textit{Paper~XVII}, 2026.
\bibitem{Kato1966} T.~Kato,
  \textit{Perturbation Theory for Linear Operators}, Springer, 1966.
\bibitem{ReedSimon4} M.~Reed, B.~Simon,
  \textit{Methods of Modern Mathematical Physics, Vol.~IV},
  Academic Press, 1978.
\bibitem{Zettl2005} A.~Zettl,
  \textit{Sturm--Liouville Theory}, AMS, 2005.
\bibitem{Olver1997} F.~Olver,
  \textit{Asymptotics and Special Functions}, AK Peters, 1997.
\bibitem{TracyWidom1994} C.~Tracy, H.~Widom,
  Level-spacing distributions and the Airy kernel,
  \textit{Comm.\ Math.\ Phys.}\ \textbf{159} (1994), 151--174.
\bibitem{Simon2005} B.~Simon,
  \textit{Trace Ideals and Their Applications}, 2nd ed., AMS, 2005.
\bibitem{DavisKahan1970} C.~Davis, W.~Kahan,
  The rotation of eigenvectors by a perturbation~III,
  \textit{SIAM J.\ Numer.\ Anal.}\ \textbf{7} (1970), 1--46.
\end{thebibliography}

\end{document}
